% Source: https://artofproblemsolving.com/community/c3h1448892_cool_triangle_fact__expand_on_it
% Stable identity: aops-topic-1448892-q01
The triangle inequality for sides $a,b,c$ requires
\[
a<b+c,\qquad b<a+c,\qquad c<a+b .
\]
Since $a\ge b$, the second is implied by the first ($c>a-b\ge0\ge b-a$). So the condition is $a-b<c<a+b$, and conversely any such $c$ gives a triangle. The integers strictly between $a-b$ and $a+b$ are $a-b+1,\dots,a+b-1$, and there are
\[
(a+b-1)-(a-b+1)+1=\boxed{2b-1}
\]
of them, where $b$ is the shorter given side. Since $b\ge1$, this is always at least $1$ (for $a=b$ it counts $c=1,\dots,2a-1$). For example, with sides $7$ and $3$ the third side can be $5,6,7,8,9$: five values, $2\cdot3-1$.
\begin{asy}
usepackage("amsmath");
size(4cm);
defaultpen(fontsize(10pt));
pair A=(0,0),B=(7,0),C=(2.2,3);
draw(A--B--C--cycle,linewidth(0.9));
label("$a$",(A+B)/2,S); label("$b$",(A+C)/2,NW); label("$c$",(B+C)/2,NE);
label("\mbox{}",(min(currentpicture,user=true)+max(currentpicture,user=true))/2,black); // keeps the text after the figure black when LaTeX typesets the labels inline
\end{asy}
